\chapter{Introduction}%
\label{cha:Introduction}
\begin{parag}{What is computer security}
    The first question that we'll answer is what are we studying in this course, what is this course about 
	\begin{definition}{Computer Security}
		\important{Properties} of a computer system must hold in the presence of a resourced strategic adversary.
	\end{definition}
\end{parag}
\begin{parag}{What Properties?}
    Traditional properties are such 
	\begin{itemize}
	    \item \important{Confidentiality} --- prevention of unauthorized disclosure of information (e.g. The adversary should not be able to read my bank statement)
	    \item \important{Integrity} --- prevention of unauthorized modification of information (e.g. The adversary should not be able to change my bank balance)
	    \item \important{Availability} --- prevention of unauthorized denial of service (e.g. The adversary should not prevent me accessing my bank account)
	\end{itemize}
	But there is also not-so-traditional (But important properties)
	\begin{itemize}
	    \item Authenticity
	    \item Anonymity
	    \item Isolation
	    \item Non-repudiation
	    \item $\vdots$
	\end{itemize}
\end{parag}

\begin{parag}{The Security policy}
	\begin{definition}{Security policy}
    A high-level description of the \important{security properties} that must hold in the system in relation to \important{assets} and \important{principals}.
    \end{definition}
	\begin{itemize}
	    \item \important{Assets} (objects): anything with value that needs to be protected
	    \item \important{Principals} (subjects): people, computer program, services
	\end{itemize}
\end{parag}
\begin{parag}{The resources strategic adversary}
    Coming back on the definition of computer security, what do we mean by '\textit{resourced strategic adversary}'? To do so we need to define the \important{threat model}
	\begin{definition}{Threat model}
	Describes the \important{resources} available to the \important{adversary} and the adversary's capabilities (observe, influence, corrupt, ...)
	\end{definition}
	\begin{itemize}
	    \item The \important{adversary} is a malicious entity aiming at breaching the security policy
	    \item The \important{adversary} is \important{strategic}: the adversary will choose the \textbf{optimal} way to use her resource to mount an attack that violates the security properties.
	\end{itemize}
	(Threat modelling is a very hard task!)
\end{parag}
Here we have defined a lot of key vocabulary to talk about the adversary but we miss a few and here are those:
\begin{definition}{Threat}
What is the feared event, the goal of the adversary that we don’t want materialized
\end{definition}
For instance 
\begin{center}
    \textit{A hacker wants to retrieve money breaking into the bank’s system\\ A student wants to learn my password by looking over my shoulder}
\end{center}

\begin{definition}{Vulnerability}
Specific weakness that could be exploited by adversaries with interest in a lot of different assets
\end{definition}

For instance 
\begin{center}
    \textit{The banking API is not protected. The password appears in plain text in my screen.}
\end{center}

\begin{definition}{Harm}
The bad thing that happens when the threat materializes
\end{definition}
For instance
\begin{center}
    \textit{The adversary steals money, he blocks access to the bank\ldots}
\end{center}

All together this give if we were to expand the computer security definition 
\begin{center}
    Properties define by a security policy of a computer system must hold in presence of a resource strategic adversary described by the threat model
\end{center}
If the \important{properties} hold within the \important{threat model}, it means that there are no \textbf{Vulnerabilities} that can be exploited to materialize threats and no harms happens.

\section{Basic concepts of security engineering}

\begin{parag}{Securing a system}
    In order to secure a system one needs a \important{security mechanism}
	\begin{definition}{Security mechanism}
	Technical mechanism used to ensure that the security policy is not violated by an adversary within the threat model
	\end{definition}
	This can be build with software (programs) + Hardware + Maths (cryptography) \& distributed systems, people and procedures. Security mechanisms can be engineered!
	\begin{subparag}{Example 1}
	    \textbf{Policy}:ensure the log of transactions is not tampered with by a single employee\\
		\textbf{Mechanism}: encrypt the message before sending
	\end{subparag}
	\begin{center}
		Systems are big! We need security mechanism\textbf{S}
	\end{center}
	Security does not necessarily increase linearly with the number of mechanisms. We could create hundred of walls with locked doors while hiding the key below the entrance carpet.
\end{parag}

\begin{parag}{How do we show systems are secure?}
    Here we can see that there is a big asymmetry between the attacker and the defender. On the attacker case, we only need \important{one way} to violate \important{one} security property (within the threat model). The defender on the other hand need that \important{no adversary strategy} can violate the security policy.\\
	$ $\\
	Therefore one can see that the question '\textit{is this system secure?}' is not enough, the correct question is 
	\begin{center}
	    Is this system secure under \important{this} threat model?
	\end{center}
	A system is said 'secure' if an adversary constrained by a \important{Specific threat model} cannot violate the \important{security policy}.\\
	$ $\\
	Then in order to show that systems are secure we need arguments:
	\begin{definition}{Security argument}
	rigorous argument that the security mechanisms in place are indeed effective in maintaining the security policy (verbal of mathematical).\\
	\begin{center}
	    Subject to the assumptions of the threat model
	\end{center}
	\end{definition}
	For a threat model to be useful, the model must constrain the adversary, otherwise we cannot make a security argument. One could just state: \textit{the adversary know everything and can do anything}.

\end{parag}


\begin{parag}{Security engineering}
    In security engineering we have three steps 
	\begin{enumerate}
	    \item High-level specification
				\begin{itemize}
				    \item Define the architecture of the system (e.g., high level block diagram)
				    \item Define the security policy (principals, assets, security properties)
					\item Define the threat model 
				\end{itemize}
		\item  Security design
			\begin{itemize}
			    \item Select/Design security mechanisms
			    \item State your security argument: which controls maintain which properties?
			\end{itemize}
		\item Secure implementation
			\begin{itemize}
			\item \important{Implement} mechanisms
			\item Ensure they \important{conform} to the design model
			\item Security testing
			\end{itemize}
	\end{enumerate}
	\begin{subparag}{Summary}
	    Security problems always involve an adversary.
		\begin{itemize}
			\item  The adversary is strategic, will take the most damaging approach
			\item The adversary’s capabilities define a threat model
			\item Security mechanisms aim at fulfilling a security policy within a threat model
			\item Showing security implies providing a security argument
		\end{itemize} 
	\end{subparag}
\end{parag}












